\subsection{NORMALLY CONVERGENT INTEGRALS}

Let \((\mathbf{f}_{\alpha})_{\alpha\in\mathrm{A}}\) be a family of regulated functions on an arbitrary interval \(I\subset R\) , with values in a complete normed space E over R. Suppose that there exists a finite real regulated function g on I such that, for every \(x\in I\) and every \(\alpha\in A\) , \(\|f_{\alpha}(x)\|\leqslant g(x)\) and also the integral \(\int_{I}g(t)dt\) is convergent. Under these conditions the integral \(\int_{I}\mathbf{f}_{\alpha}(t)dt\) is absolutely and uniformly convergent on A; in fact, for every compact interval K contained in I,

\[
\left\| \int_ {K} \mathbf {f} _ {\alpha} (t) d t \right\| \leqslant \int_ {K} g (t) d t
\]

and the convergence of the integral \(\int_{I} g(t) dt\) implies that for every \(\varepsilon > 0\) there exists a compact interval \(J \subset I\) such that for every compact interval \(K \subset I\) disjoint from J one has \(\int_{K} g(t) dt \leqslant \varepsilon\) . When there exists a real function g having the preceding properties one says that the integral \(\int_{I} f_{\alpha}(t) dt\) is normally convergent on A (cf.~Gen.~Top., X, p.~296).

An integral can be uniformly convergent on A without being normally convergent. *This happens for the sequence \((f_n)\) of real functions defined by the conditions \(f_n(x) = 1 / x\) for \(n \leqslant x \leqslant n + 1\) , and \(f_n(x) = 0\) for the other values of \(x\) in \(\mathbf{I} = [\mathbf{0}, +\infty[\) . It is immediate that the integral \(\int_{1}^{\infty} f_n(t) dt\) is uniformly convergent, but not normally convergent, since the relation \(g(x) \geqslant f_n(x)\) for each \(x \in \mathbf{I}\) and all \(n\) entails that \(g(x) \geqslant 1 / x\) , and consequently that the integral of \(g\) over \(\mathbf{I}\) is not convergent.*

In particular, let us consider a series whose general term \(u_{n}\) is a regulated function on an interval I, and suppose that the series with general term \(\|u_{n}(x)\|\) (which is a regulated function on I) converges uniformly on every compact interval contained in I, and such that the series with general term \(\int_{I}\|\mathbf{u}_{n}(t)\|dt\) is convergent; then (II, p.~66, prop. 2) the (regulated) function \(g(x)\) , the sum of the series with general term \(\|u_{n}(x)\|\) , is such that the integral \(\int_{I}g(t)dt\) is convergent. If one puts \(f_{n}=\sum_{p=1}^{n}u_{p}\) , then the integral \(\int_{I}f_{n}(t)dt\) is normally convergent, for one has

\[
\left\| \mathbf {f} _ {n} (x) \right\| \leqslant \sum_ {p = 1} ^ {n} \left\| \mathbf {u} _ {p} (x) \right\| \leqslant g (x)
\]

for all \(x \in \mathrm{I}\) and all \(n\) ; in consequence, the sum \(\mathbf{f}\) of the series with general term \(\mathbf{u}_n\) is a regulated function on \(\mathrm{I}\) such that the integral \(\int_{\mathrm{I}} \mathbf{f}(t) dt\) is convergent, and one has

\[
\int_ {\mathrm{I}} \mathbf {f} (t) d t = \sum_ {n = 1} ^ {\infty} \int_ {\mathrm{I}} \mathbf {u} _ {n} (t) d t\tag{11}
\]

(``term-by-term integration of a series on a non-compact interval'').

